\setcounter{equation}{0}
\setcounter{Q}{0}

Les alliages de titane et de vanadium sont utilisés dans le secteur aéronautique, pour la réalisation des réacteurs et des trains d’atterrissage. Le diagramme binaire isobare solide-liquide simplifié et limité aux hautes températures est représenté à la figure suivante sous une pression $P = P^\circ = 1 ~ \si{\bar}$, avec en abscisse la fraction massique en vanadium, $w(\ce{V})$, et en ordonnée la température exprimée en \si{\celsius}.

\begin{center}
\begin{minipage}{0.48\textwidth}
\includegraphics[width=\textwidth]{/home/jeff/Documents/IUTOrsay/1-Cours/4p13-Alliages/Cours-TD/Exos/Binaires-Lecture/Homogene-Ex04/figures/Fig01}
\end{minipage}\begin{minipage}{0.48\textwidth}
\captionof{figure}{Diagramme binaire isobare, à une pression de \SI{1}{\bar}, du mélange \ce{Ti}-\ce{V}.}
[La composition est donnée par la fraction massique en vanadium dans le mélange]
\end{minipage}
\end{center}

\Q Indiquer le nombre et la nature des phases en présence dans les domaines \textbf{A}, \textbf{B} et \textbf{C}.

\Q Un point remarquable \textbf{I} apparaît sur ce diagramme binaire pour une fraction massique en vanadium $w = 0.3$ et une température $\theta = 1560 ~ \si{\celsius}$. Préciser la propriété physique remarquable du mélange correspondant.

\Q Représenter l’allure des courbes d’analyse thermique isobare de refroidissement pour des fractions massiques en vanadium respectivement égales à $1.0$ ; $0.1$ et $0.3$. Justifier votre réponse par un calcul de variance pour cette dernière courbe d’analyse thermique.

\Q Lors du refroidissement, à partir de $\theta = 1750 ~ \si{\celsius}$, du mélange représenté sur le diagramme par le point ($\bullet$), donner la température d’apparition du premier cristal de solide et déterminer la composition en fraction massique de ce premier cristal de solide.

\Q Un mélange liquide titane-vanadium est préparé à partir de \SI{100}{\kilo\gram} de vanadium et de \SI{900}{\kilo\gram} de titane. Ce mélange est porté à $\theta = 1600 ~ \si{\celsius}$. Indiquer la nature et la composition en fraction massique des phases en équilibre à cette température. Calculer les masses de titane et de vanadium dans chacune des phases.